\documentclass[tikz,border=6pt]{standalone}
\usepackage{ctex}
\usepackage{amsmath}
\usepackage{pgfplots}
\pgfplotsset{compat=1.18}
\usetikzlibrary{calc,angles,quotes,arrows.meta,positioning,decorations.markings}
\begin{document}
\begin{tikzpicture}[>=Stealth]

  % ===== 参数 =====
  \def\A{1.6}        % 振幅（旋转向量长度）
  \def\phi{30}       % 当前相位角（度）
  \pgfmathsetmacro\px{\A*cos(\phi)}
  \pgfmathsetmacro\py{\A*sin(\phi)}

  % ====================================================
  % 左侧：复平面上的旋转相量（单位圆）
  % ====================================================
  \begin{scope}
    % 坐标轴
    \draw[->,gray!70] (-2.1,0) -- (2.2,0) node[right,black]{\small 实轴 Re};
    \draw[->,gray!70] (0,-2.1) -- (0,2.2) node[above,black]{\small 虚轴 Im};
    % 单位圆（半径 = 振幅 A）
    \draw[gray!55,thick] (0,0) circle (\A);
    % 相位角弧
    \draw[->,red!70!black,thick] (0.55,0) arc (0:\phi:0.55);
    \node[red!70!black] at (\phi/2:0.85) {\small $\omega t+\varphi$};
    % 旋转相量
    \draw[->,blue!70!black,very thick] (0,0) -- (\px,\py)
       node[midway,above left=-2pt,blue!70!black]{\small $A$};
    \fill[blue!70!black] (\px,\py) circle (1.6pt);
    % 旋转方向箭头
    \draw[->,orange!85!black,thick]
       (130:\A+0.28) arc (130:175:\A+0.28);
    \node[orange!85!black] at (152:\A+0.62) {\small $\omega$};
    % 实轴投影（虚线）→ 这是余弦信号的取值
    \draw[densely dashed,green!45!black,thick] (\px,\py) -- (\px,0);
    \fill[green!45!black] (\px,0) circle (1.6pt);
    \draw[->,green!45!black,thick] (0,0) -- (\px,0)
       node[midway,below=2pt,green!45!black]{\small $A\cos(\omega t+\varphi)$};
    % 标题
    \node[align=center] at (0,-2.7) {\small 左：以角速度 $\omega$ 旋转的相量};
  \end{scope}

  % ====================================================
  % 右侧：对应的正弦（余弦）波 —— 用 pgfplots
  % 横轴向右为时间，纵轴为相量在实轴上的投影
  % ====================================================
  \begin{scope}[xshift=4.3cm,yshift=0cm]
    \begin{axis}[
        width=8.4cm, height=6.0cm,
        at={(0,0)}, anchor=west,
        axis lines=middle,
        xmin=0, xmax=7.2, ymin=-2.0, ymax=2.0,
        xtick=\empty, ytick={-1.6,1.6},
        yticklabels={$-A$,$A$},
        xlabel={\small $t$（时间）}, ylabel={\small 信号},
        every axis x label/.style={at={(ticklabel* cs:1.0)},anchor=west},
        every axis y label/.style={at={(ticklabel* cs:1.0)},anchor=south},
        clip=false,
      ]
      % 余弦曲线 A cos(omega t + phi)，phi = 30deg = pi/6
      % 这里取 omega=1，相位 pi/6
      \addplot[domain=0:7,samples=220,blue!70!black,very thick]
        {1.6*cos(deg(x) + 30)};
      % 当前时刻 t0 对应左图相位：cos(omega t0 + phi)=cos(phi) → t0=0
      % 用一条水平虚线连接左图投影点高度
      \pgfmathsetmacro\yval{1.6*cos(30)}
      \draw[densely dashed,green!45!black,thick]
        (axis cs:-4.3,\yval) -- (axis cs:0,\yval);
      \fill[green!45!black] (axis cs:0,\yval) circle (1.8pt);
      \draw[densely dashed,green!45!black]
        (axis cs:0,\yval) -- (axis cs:0,0);
      % 标注一个周期
      \draw[<->,gray!70] (axis cs:0.0,-1.85) -- (axis cs:6.283,-1.85);
      \node[gray!50!black,fill=white,fill opacity=0.7,text opacity=1,inner sep=1pt]
        at (axis cs:3.14,-1.85) {\small 周期 $T=\dfrac{2\pi}{\omega}$};
    \end{axis}
    \node[align=center] at (4.0,-2.7) {\small 右：投影随时间画出 $A\cos(\omega t+\varphi)$};
  \end{scope}

\end{tikzpicture}
\end{document}
